\documentclass[11pt]{article}
\usepackage[margin=1in]{geometry}
\usepackage{amsmath,amssymb,amsthm}
\usepackage[hidelinks]{hyperref}

\newtheorem*{theorem}{Theorem}

\title{Literature Status for Question Q:UC of arXiv:1510.00276}
\author{FA/Banach run packet}
\date{June 28, 2026}

\begin{document}
\maketitle

\section*{Status}

\textbf{Result type:} literature already answered.

\textbf{Source question:} Hytönen--Li--Naor,
\emph{Quantitative affine approximation for UMD targets}, arXiv:1510.00276,
Question Q:UC in Section 1, ``Open questions.''

\textbf{Supporting answer:} Hytönen--Naor,
\emph{Heat flow and quantitative differentiation}, arXiv:1608.01915,
Theorem 1.

\section*{Source Question}

The source paper asks whether the improved lower bound for the modulus of
uniform affine approximability, proved there for UMD targets, remains valid
when the target is merely uniformly convex. In its notation, Question Q:UC asks
for an improved bound as in Theorem 1 under the weaker assumption that \(Y\)
admits an equivalent uniformly convex norm, and notes that this is the maximal
qualitative generality by the theorem of Bates--Johnson--Lindenstrauss--Preiss--
Schechtman.

\section*{Supporting Literature Answer}

The later paper arXiv:1608.01915 states the following theorem in its
introduction.

\begin{theorem}
Suppose that \(Y\) admits an equivalent uniformly convex norm. Then there is
\(c=c(Y)\in(0,\infty)\) such that for every \(n\in\mathbb N\), every
\(n\)-dimensional normed space \(X\), and every \(\varepsilon\in(0,1/2]\),
\[
  r^{X\to Y}(\varepsilon)\ge \exp\left(-\frac{1}{\varepsilon^{c n}}\right).
\]
\end{theorem}

Immediately after the theorem, the later paper states that this result answers
Question 8 of arXiv:1510.00276 positively. This is a direct supporting-paper
identification rather than an agent-inferred implication.

\section*{Scope and Limitations}

This packet records the positive literature answer to the uniformly
convex-target part of Question Q:UC. It does not settle the broader problem of
the sharp asymptotic rate of \(r_n(\varepsilon)\) in the Hilbert-to-Hilbert
case, nor the separate infinite-dimensional-domain question from the same
source section.

The supporting paper also improves the Hilbertian \(L_2\) Dorronsoro estimate,
but it still records the quantitative differentiation asymptotic problem as
open in the sense that the best known Hilbert-to-Hilbert bound for
\(r^{\ell_2^n\to\ell_2}(1/4)\) has the same asymptotic form as the general
bound.

\section*{Search Bounds}

Before packaging, the run's registry, solutions index, attempts index, and
proof-gaps index were searched for arXiv:1510.00276, the paper title,
``affine approximation,'' ``quantitative differentiation,'' and ``UMD targets.''
No prior packet for arXiv:1510.00276 was found. The nearby arXiv:1608.01915
registry hit concerns Question 19 / Q:proj lip of the later paper, not this
source question.

\section*{References}

\begin{thebibliography}{9}

\bibitem{HLN2017}
T. Hytönen, S. Li, and A. Naor,
\emph{Quantitative affine approximation for UMD targets},
Discrete Analysis 2016:37, 1--71.
arXiv:1510.00276.
\url{https://arxiv.org/abs/1510.00276}

\bibitem{HN2016}
T. Hytönen and A. Naor,
\emph{Heat flow and quantitative differentiation},
J. Eur. Math. Soc. 20 (2018), no. 9, 2183--2247.
arXiv:1608.01915.
\url{https://arxiv.org/abs/1608.01915}

\end{thebibliography}

\end{document}
